\begin{solucion}
    Queremos contar las permutaciones de $[n]=\{1,2,\dots,n\}$ en las que
\emph{ningún número par queda fijo}.  
Sea $E=\{\,2,4,\dots,2\lfloor n/2\rfloor\}$ el conjunto de los $m=\lfloor n/2\rfloor$
índices pares.

\textbf{Paso clave (Principio de Inclusión–Exclusión).}  
Para cada $e\in E$ defina 
\[
A_e=\{\sigma\in S_n:\,\sigma(e)=e\},
\]
el evento “$e$ es punto fijo”.
Nuestro objetivo es 
\(\bigl|\bigcap_{e\in E}\overline{A_e}\bigr|\).

Si $J\subseteq E$ contiene $j$ índices pares, fijar \emph{todos} los puntos de $J$
obliga a dejar invariables precisamente esos $j$ elementos y
permite permutar libremente los $n-j$ restantes, de modo que  
\[
|A_J|=\bigl|\,\bigcap_{e\in J}A_e\bigr|=(n-j)! .
\]
Además hay $\binom{m}{j}$ posibles subconjuntos $J$ de tamaño $j$.

Aplicando inclusión–exclusión obtenemos
\[
\#\{\sigma\in S_n:\, \text{ningún par fijo}\}
=\sum_{j=0}^{m}(-1)^{\,j}\binom{m}{j}\,(n-j)!.
\]

\textbf{Resultado final.}\quad
Si $m=\lfloor n/2\rfloor$, el número buscado es
\[
\boxed{\displaystyle
\sum_{j=0}^{m}(-1)^{\,j}\binom{m}{j}\,(n-j)!
}.
\]

Para $n\le 8$ los primeros valores son  
\(
1,\;1,\;4,\;14,\;78,\;426,\;3216,\;24024
\),
coincidiendo con la fórmula.
\end{solucion}\bigskip
